%  DESARROLLO PASO A PASO — EJERCICIO 1.3
%  Inecuación con raíz cuadrada: x + 1 < sqrt(x^2 + 3x)
%  Usach Premium — Cálculo I
%  Compila con: pdflatex c1_desarrollo_1_3.tex  (2 veces)
% ============================================================

\documentclass[letterpaper, 12pt]{article}

% ============================================================
% PAQUETES
% ============================================================
\usepackage[utf8]{inputenc}
\usepackage[spanish]{babel}
\usepackage{amsmath, amssymb}
\usepackage{geometry}
\usepackage{fancyhdr}
\usepackage{enumitem}
\usepackage{graphicx}
\usepackage{hyperref}
\usepackage{parskip}
\usepackage{xcolor}
\usepackage{tcolorbox}
\usepackage{eso-pic}
\usepackage{tikz}
\usepackage{array}
\usepackage{booktabs}
\usepackage{colortbl}
\usetikzlibrary{patterns, arrows.meta, decorations.pathmorphing, babel}
\tcbuselibrary{skins, breakable}

% ============================================================
% GEOMETRÍA
% ============================================================
\geometry{letterpaper, margin=2cm, top=3cm, bottom=2.5cm, headheight=2cm}

% ============================================================
% COLORES INSTITUCIONALES
% ============================================================
\definecolor{celesteSuave}{HTML}{F0F7FF}
\definecolor{azulFuerte}{HTML}{1A5276}
\definecolor{verdePaso}{HTML}{1E8449}
\definecolor{rojoAlerta}{HTML}{C0392B}
\definecolor{naranjaAcento}{HTML}{D35400}
\definecolor{grisTabla}{HTML}{D5D8DC}
\definecolor{enlace}{HTML}{1A5276}

% ============================================================
% RUTAS DE IMÁGENES
% ============================================================
\newcommand{\logoup}{../../../../../template/images/logo_up.png}
\newcommand{\logousach}{../../../../../template/images/logo_usach.png}
\newcommand{\logofondo}{../../../../../template/images/logo_up.png}

% ============================================================
% INFORMACIÓN
% ============================================================
\newcommand{\institucion}{Usach Premium}
\newcommand{\curso}{Cálculo I}
\newcommand{\guiasnumero}{Desarrollo — Ejercicio 1.3}
\newcommand{\semestre}{1er. Semestre de 2026}
\newcommand{\preparadopor}{Usach Premium.}
\newcommand{\webinstitucion}{www.usachpremium.cl}
\newcommand{\instagraminstitucion}{https://www.instagram.com/usach.premium}
\newcommand{\transparencia}{0.04}
\newcommand{\fraseportada}{"Por y para estudiantes"}

% ============================================================
% HYPERREF
% ============================================================
\hypersetup{colorlinks=true, linkcolor=enlace, urlcolor=enlace,
            citecolor=enlace, pdfborder={0 0 0}}

% ============================================================
% ENCABEZADO Y PIE
% ============================================================
\pagestyle{fancy}
\fancyhf{}
\fancyhead[L]{\includegraphics[height=1.2cm]{\logoup}}
\fancyhead[R]{%
    \begin{minipage}[b]{0.45\textwidth}
        \raggedleft\fontsize{9pt}{11pt}\selectfont
        \textbf{\color{azulFuerte}\institucion} \\
        \curso\ --- \guiasnumero \\
        \textit{\color{gray}@usach.premium}
    \end{minipage}\hspace{0.1cm}%
    \includegraphics[height=1.2cm]{\logousach}}
\fancyfoot[L]{\fontsize{9pt}{11pt}\selectfont \textit{\fraseportada}}
\fancyfoot[R]{\fontsize{9pt}{11pt}\selectfont P\'agina \thepage}
\renewcommand{\headrulewidth}{0.4pt}
\renewcommand{\footrulewidth}{0.4pt}
\fancypagestyle{plain}{\fancyhf{}\renewcommand{\headrulewidth}{0pt}
                                   \renewcommand{\footrulewidth}{0pt}}

% ============================================================
% MARCA DE AGUA
% ============================================================
\newcommand{\MarcaAgua}{
    \AddToShipoutPicture{
        \AtPageCenter{
            \begin{tikzpicture}[remember picture, overlay]
                \node[opacity=\transparencia] at (0,0)
                    {\includegraphics[width=0.7\textwidth]{\logofondo}};
            \end{tikzpicture}}}}

% ============================================================
% CAJAS
% ============================================================
% Caja de paso numerado
\newtcolorbox{paso}[2][]{
    colback=celesteSuave, colframe=azulFuerte, fonttitle=\bfseries,
    coltitle=white, colbacktitle=azulFuerte,
    title={\large Paso #2},
    arc=6pt, boxrule=1pt, enhanced, #1
}

% Caja de advertencia/justificación
\newtcolorbox{justif}[1]{
    colback=naranjaAcento!8, colframe=naranjaAcento, fonttitle=\bfseries\small,
    title={#1}, arc=5pt, boxrule=0.8pt, enhanced
}

% ============================================================
\begin{document}
\thispagestyle{plain}
\MarcaAgua

% ============================================================
% PORTADA
% ============================================================
\AddToShipoutPicture*{%
    \put(54,700){\includegraphics[width=2cm]{\logoup}}
    \put(420,706){\includegraphics[width=5cm]{\logousach}}}

\vspace*{0.5cm}
\begin{center}
    \color{azulFuerte}
    {\huge \textbf{DESARROLLO PASO A PASO}} \\[0.5cm]
    {\Huge \textbf{\curso}} \\[0.4cm]
    {\Large Ejercicio 1.3 --- Inecuaci\'on con ra\'iz cuadrada} \\[1.2cm]
    \begin{tcolorbox}[colback=celesteSuave, colframe=azulFuerte, arc=10pt,
        width=0.88\textwidth, center, boxrule=1.5pt]
        \centering\vspace{0.3cm}
        {\large \textbf{Inecuaci\'on a resolver}}\\[0.6cm]
        \color{black}
        {\LARGE $\displaystyle x + 1 \;<\; \sqrt{x^2 + 3x}$}
        \vspace{0.4cm}
    \end{tcolorbox}
\end{center}

\vspace{1.2cm}
\begin{center}
    {\Large \textbf{\semestre}} \\[0.7cm]
    {\large \textbf{\preparadopor}}
\end{center}

\vspace{1.8cm}
\begin{center}
    {\color{azulFuerte}\hrule height 0.5pt}\vspace{0.3cm}
    \begin{minipage}{0.45\textwidth}\centering
        \textbf{Instagram:} \\\small\href{\instagraminstitucion}{@usach.premium}
    \end{minipage}
    \begin{minipage}{0.45\textwidth}\centering
        \textbf{Web:} \\\small\href{http://\webinstitucion}{\webinstitucion}
    \end{minipage}
\end{center}

\pagenumbering{arabic}
\newpage

% ============================================================
% ENUNCIADO
% ============================================================
\begin{tcolorbox}[colback=azulFuerte!12, colframe=azulFuerte, arc=8pt,
    boxrule=1.2pt, enhanced]
    \textbf{\large Ejercicio 1.3.} \quad Resuelva la siguiente inecuaci\'on:
    \[
        x + 1 \;<\; \sqrt{x^2 + 3x}
    \]
\end{tcolorbox}

\bigskip

% ============================================================
% PASO 1 — DOMINIO
% ============================================================
\begin{paso}{1: Determinaci\'on del Dominio}

Para que la expresi\'on est\'e definida en $\mathbb{R}$, el radicando debe ser
no negativo:
\[
    x^2 + 3x \;\geq\; 0
\]

Factorizando:
\[
    x(x + 3) \;\geq\; 0
\]

\textbf{Puntos cr\'iticos:} $x = 0$ y $x = -3$.

\medskip
\begin{center}
\renewcommand{\arraystretch}{1.5}
\begin{tabular}{c|ccccc}
    & $x < -3$ & $x = -3$ & $-3 < x < 0$ & $x = 0$ & $x > 0$ \\
\hline
$x$       & $-$ & $-$ & $-$ & $0$ & $+$ \\
$(x + 3)$ & $-$ & $0$ & $+$ & $+$ & $+$ \\
\hline
\rowcolor{celesteSuave}
$x(x+3)$  & $+$ & $0$ & $-$ & $0$ & $+$ \\
\end{tabular}
\end{center}

\medskip
$x(x+3) \geq 0$ cuando $x \leq -3$ \;\'o\; $x \geq 0$.

\[
    \boxed{Dom\,f \;=\; (-\infty,\,-3\,] \;\cup\; [\,0,\,+\infty)}
\]

\end{paso}

\bigskip

% ============================================================
% PASO 2 — ANÁLISIS POR CASOS
% ============================================================
\begin{paso}{2: An\'alisis por casos seg\'un el signo de $x + 1$}

El lado derecho $\sqrt{x^2+3x} \geq 0$ siempre (dentro del dominio).
El lado izquierdo $x + 1$ puede ser negativo, cero o positivo.
Analizamos por separado seg\'un si el LHS es negativo o no negativo,
pues \textbf{solo podemos elevar al cuadrado cuando ambos lados son no negativos}.

\bigskip
\textbf{Caso 1: $x + 1 < 0$, es decir, $x < -1$}

\medskip
En el dominio, $x < -1$ implica $x \leq -3$.
Como $x + 1 < 0 \leq \sqrt{x^2+3x}$, la desigualdad se satisface
\textbf{autom\'aticamente} para todo $x \leq -3$.

\begin{justif}{Justificaci\'on}
    Un n\'umero negativo es siempre menor que un n\'umero no negativo.
    Si $A < 0$ y $B \geq 0$, entonces $A < B$ sin condici\'on alguna.
\end{justif}

\[
    \text{Soluci\'on parcial (Caso 1): } x \in (-\infty,\,-3\,]
\]

\bigskip
\textbf{Caso 2: $x + 1 \geq 0$, es decir, $x \geq -1$}

\medskip
En el dominio, $x \geq -1$ implica $x \geq 0$ (pues el dominio excluye
$(-3, 0)$).
Ahora \textbf{ambos lados son no negativos}, por lo que podemos elevar al
cuadrado preservando la desigualdad.

\end{paso}

\bigskip

% ============================================================
% PASO 3 — ELEVACIÓN AL CUADRADO (CASO 2)
% ============================================================
\begin{paso}{3: Elevaci\'on al cuadrado --- Caso 2 ($x \geq 0$)}

\begin{justif}{¿Por qu\'e podemos elevar al cuadrado?}
    Si $0 \leq A < B$, entonces $A^2 < B^2$.
    Aqu\'i $A = x+1 \geq 0$ y $B = \sqrt{x^2+3x} \geq 0$,
    luego la equivalencia es v\'alida y \textbf{no cambia el sentido
    de la desigualdad}.
\end{justif}

\[
    x + 1 \;<\; \sqrt{x^2 + 3x}
    \quad\stackrel{(\,)^2}{\Longleftrightarrow}\quad
    (x+1)^2 \;<\; x^2 + 3x
\]

Expandiendo el lado izquierdo:
\[
    x^2 + 2x + 1 \;<\; x^2 + 3x
\]

Simplificando (restamos $x^2$ en ambos lados):
\[
    2x + 1 \;<\; 3x
\]

Reordenando:
\[
    1 \;<\; x \qquad\Longleftrightarrow\qquad x \;>\; 1
\]

\[
    \text{Soluci\'on parcial (Caso 2): } x \in (1,\,+\infty) \;\text{ con } x \geq 0
\]

\end{paso}

\newpage

% ============================================================
% PASO 4 — TABLA DE SIGNOS (expresión simplificada)
% ============================================================
\begin{paso}{4: Tabla de signos de $p(x) = (x+1)^2 - (x^2+3x)$}

Tras elevar al cuadrado en el Caso 2, la inecuaci\'on se redujo a:
\[
    (x+1)^2 - (x^2+3x) \;<\; 0
\]

Calculamos expl\'icitamente esta diferencia:
\[
    (x+1)^2 - (x^2+3x)
    \;=\; x^2 + 2x + 1 - x^2 - 3x
    \;=\; 1 - x
\]

\textbf{Punto cr\'itico:} $p(x) = 1 - x = 0 \;\Rightarrow\; x = 1$.

\medskip
\begin{center}
\renewcommand{\arraystretch}{1.7}
\begin{tabular}{c|ccc}
    & $x < 1$ & $x = 1$ & $x > 1$ \\
\hline
$p(x) = 1 - x$ & $+$ & $0$ & $-$ \\
\hline
\rowcolor{celesteSuave}
Signo & \color{rojoAlerta}$p > 0$ & $p = 0$ & \color{verdePaso}$\mathbf{p < 0}$ \\
\end{tabular}
\end{center}

\medskip
La condici\'on $p(x) < 0$ se cumple exactamente cuando $x > 1$,
confirmando el resultado del Paso 3.

\bigskip
\textbf{Recta num\'erica — Caso 2:}

\begin{center}
\begin{tikzpicture}[scale=1.2]
    \draw[->] (-0.5,0) -- (4.5,0) node[right] {$x$};
    % x=0 (borde dominio)
    \draw[azulFuerte!60, dashed, thick] (0,-0.55) -- (0,0.55);
    \fill[azulFuerte!60] (0,0) circle (2pt);
    \node[azulFuerte] at (0,-0.55) {\small$0$};
    % x=1 (punto critico, excluido)
    \draw[thick] (1,0) circle (3.5pt);
    \node[below=5pt] at (1,0) {\small$1$};
    % signos
    \node at (0.5, 0.4) {\color{rojoAlerta}\small$p > 0$};
    \node at (2.8, 0.4) {\color{verdePaso}\small$p < 0$};
    % segmento solución
    \draw[verdePaso, very thick, ->] (1.07,0) -- (4.4,0);
    \node[verdePaso] at (3.0,-0.55) {\small Soluci\'on Caso 2: $(1,+\infty)$};
\end{tikzpicture}
\end{center}

\end{paso}

\bigskip

% ============================================================
% PASO 5 — UNIÓN DE CASOS E INTERSECCIÓN CON DOMINIO
% ============================================================
\begin{paso}{5: Uni\'on de casos e intersecci\'on con el dominio}

Reunimos los resultados de ambos casos:

\begin{itemize}[leftmargin=1.8cm, itemsep=4pt]
    \item \textbf{Caso 1} ($x+1 < 0$ en dominio): \quad
          $x \in (-\infty,\,-3\,]$
    \item \textbf{Caso 2} ($x+1 \geq 0$ en dominio): \quad
          $x \in (1,\,+\infty)$
\end{itemize}

\medskip
La soluci\'on total es la \textbf{uni\'on} de ambos conjuntos:
\[
    \mathcal{S} \;=\; (-\infty,\,-3\,] \;\cup\; (1,\,+\infty)
\]

Verificamos que esta soluci\'on est\'a contenida en el dominio:
\begin{itemize}[leftmargin=1.8cm]
    \item $(-\infty,-3] \;\subset\; (-\infty,-3] \;\subset\; Dom\,f$ \checkmark
    \item $(1,+\infty) \;\subset\; [0,+\infty) \;\subset\; Dom\,f$ \checkmark
\end{itemize}

\bigskip
\textbf{Recta num\'erica final:}

\begin{center}
\begin{tikzpicture}[scale=0.9]
    \draw[->] (-6.5,0) -- (5.0,0) node[right] {$x$};
    % x=-3 incluido
    \fill[verdePaso] (-3,0) circle (3pt);
    \node[below=5pt] at (-3,0) {\small$-3$};
    % x=0 borde dominio (no en solucion en ese tramo)
    \fill[white] (0,0) circle (3pt);
    \draw[azulFuerte!60, thick] (0,0) circle (3pt);
    \node[below=5pt] at (0,0) {\small$0$};
    % x=1 excluido
    \draw[verdePaso, thick] (1,0) circle (3.5pt);
    \node[below=5pt] at (1,0) {\small$1$};
    % segmento izquierdo solución (-inf, -3]
    \draw[verdePaso, very thick, <-] (-6.4,0) -- (-3.0,0);
    % segmento derecho solución (1, +inf)
    \draw[verdePaso, very thick, ->] (1.07,0) -- (4.8,0);
    % etiquetas de signo
    \node at (-5.0, 0.5) {\color{verdePaso}\small Satisface};
    \node at (-1.5, 0.5) {\color{rojoAlerta}\small No satisface};
    \node at ( 0.5, 0.5) {\color{rojoAlerta}\small No};
    \node at ( 3.0, 0.5) {\color{verdePaso}\small Satisface};
    % dominio (regiones)
    \node[azulFuerte] at (-3.0,-0.9) {\small Dom};
    \node[azulFuerte] at ( 2.5,-0.9) {\small Dom};
\end{tikzpicture}
\end{center}

\end{paso}

\bigskip

% ============================================================
% RESULTADO FINAL
% ============================================================
\begin{tcolorbox}[colback=azulFuerte, colframe=azulFuerte,
    colupper=white, arc=10pt, boxrule=1.5pt, enhanced,
    drop shadow={black!15}]
    \centering
    \large\textbf{RESULTADO FINAL}\\[0.5em]
    \LARGE
    $\displaystyle
        x \;\in\; (-\infty,\,-3\,] \;\cup\; (1,\,+\infty)$
\end{tcolorbox}

\vspace{0.8cm}

% ============================================================
% VERIFICACIÓN (sólo resumen, sin código Python)
% ============================================================
\begin{tcolorbox}[colback=verdePaso!8, colframe=verdePaso, arc=8pt,
    boxrule=1pt, enhanced,
    title={\bfseries\color{white} Verificaci\'on num\'erica (SymPy + NumPy)},
    colbacktitle=verdePaso, fonttitle=\bfseries]
    \small
    Se valid\'o la soluci\'on con Python usando dos librer\'ias independientes:
    \begin{itemize}[leftmargin=1.2cm, itemsep=2pt]
        \item \textbf{SymPy} (c\'alculo simb\'olico): \;
              \texttt{solveset} retorn\'o exactamente
              $(-\infty,-3]\cup(1,+\infty)$.
        \item \textbf{NumPy} (muestreo 50\,000 puntos con desfase irracional):
              acuerdo del \textbf{100.0000\,\%} entre la inecuaci\'on original
              y el conjunto soluci\'on.
    \end{itemize}
    Todos los puntos de borde verificados:
    $x = -3$ satisface (\checkmark),
    $x = 0$ no satisface (\checkmark),
    $x = 1$ no satisface (\checkmark),
    $x = 1.0001$ satisface (\checkmark).
\end{tcolorbox}

\vspace{1cm}
\begin{center}
    \small\textbf{Usach Premium} --- ¡Nos vemos en la ayudant\'ia! \\
    \textit{"Por y para estudiantes."}
\end{center}

\end{document}
